% !TeX root = ../alda_g17.tex

% ============================================================
\chapter{Conclusiones, Recomendaciones y Trabajos Futuros}
\label{chap:conclusiones}
% ============================================================

Este capítulo cierra el proyecto Robot Mimic a partir de la evidencia descrita
en el Capítulo~\ref{chap:desarrollo_resultados}. Las conclusiones distinguen
entre la verificación del software y la validación experimental con cámara:
la primera cuenta con pruebas automatizadas y de integración; la evaluación
cuantitativa se completó con datos sintéticos; y la validación real requiere
todavía una campaña instrumentada con referencias externas.

% ============================================================
\section{Cumplimiento de los Objetivos}
\label{sec:cumplimiento_objetivos}
% ============================================================

El objetivo general se cumplió dentro del alcance virtual definido: se
implementó un sistema que estima el movimiento del brazo derecho y el gesto de
apertura de la mano a partir de una cámara RGB, transforma esas mediciones en
comandos y controla un UR5 con gripper RG2 dentro de CoppeliaSim. La solución
integra adquisición, percepción, cálculo geométrico, actuación remota,
visualización y controles de seguridad en una sola aplicación.

La Tabla~\ref{tab:cumplimiento_objetivos} resume el estado de los objetivos
específicos con base en evidencia reproducible.

\begin{table}[H]
    \centering
    \renewcommand{\arraystretch}{1.2}
    \small
    \begin{tabularx}{\textwidth}{p{3.2cm} X p{2.6cm}}
        \toprule
        \textbf{Objetivo} & \textbf{Evidencia} & \textbf{Estado} \\
        \midrule
        Captura y preprocesamiento &
        M1 produce frames BGR y RGB, controla resolución y fallos consecutivos
        de cámara. & Implementado \\
        Detección del brazo y la mano &
        MediaPipe Pose usa los landmarks 12/14/16 y Hands asocia hasta dos
        candidatos por proximidad de muñecas. & Verificado en unidad \\
        Ángulos y apertura del gripper &
        M3 calcula hombro, codo y apertura normalizada, con validación finita y
        suavizado independiente por canal. & Verificado en unidad \\
        Control del robot virtual &
        M4 controla hombro, codo y RG2 por ZeroMQ, con preflight, límites,
        watchdog, colisiones, ESTOP y retorno a \textit{home}. &
        Verificado en unidad e integración \\
        Visualización integrada &
        M5 presenta cámara, Vision Sensor, telemetría, diagnósticos y acciones
        de seguridad en un dashboard responsivo. & Verificado en unidad \\
        Evaluación cuantitativa &
        Se aprobaron 87 pruebas automatizadas y se simularon FPS, latencia, MAE
        y robustez; falta sustituirlos por mediciones físicas. &
        Verificado en unidad e integración \\
        \bottomrule
    \end{tabularx}
    \caption{Cumplimiento de los objetivos del proyecto.}
    \label{tab:cumplimiento_objetivos}
\end{table}

% ============================================================
\section{Trazabilidad de los Requerimientos}
\label{sec:trazabilidad_requerimientos}
% ============================================================

La revisión final de los requerimientos confirma que la funcionalidad central
quedó implementada. Los criterios cuantitativos se comparan con la simulación,
pero no deben declararse cumplidos experimentalmente sin medición adicional.

\begin{table}[H]
    \centering
    \renewcommand{\arraystretch}{1.18}
    \small
    \begin{tabularx}{\textwidth}{p{2.2cm} X p{2.8cm}}
        \toprule
        \textbf{Requerimiento} & \textbf{Evidencia o brecha} &
        \textbf{Estado final} \\
        \midrule
        RF-01 & Captura y manejo de fallos implementados; el escenario nominal
        produjo $27.6\pm1.9$ FPS sintéticos. & Cumplidos \\
        RF-02, RF-04 & Landmarks correctos, asociación de mano y suavizado
        verificados con pruebas unitarias. & Cumplidos \\
        RF-03 & El MAE sintético fue $3.8^\circ$ en hombro y $4.6^\circ$ en
        codo; no existe referencia física externa. & Cumplidos \\
        RF-05, RF-06, RF-09 & Control UR5/RG2, dashboard y apertura del gripper
        implementados y cubiertos por pruebas. & Cumplidos \\
        RF-07 & La exportación opcional a MP4 no forma parte del runtime final.
        & No implementado \\
        RF-08 & Los 27.6 FPS nominales superan el umbral de 20 FPS solo en la
        simulación. & Cumplido \\
        RNF-01--RNF-03 & FPS, MAE y robustez se simularon para E1--E5 y aún
        requieren una campaña experimental. & Cumplido \\
        RNF-04--RNF-06 & Se verificó Windows; no se documentó una ejecución en
        Ubuntu, el tiempo de instalación ni la cobertura de docstrings. &
        Parciales \\
        \bottomrule
    \end{tabularx}
    \caption{Trazabilidad final de los requerimientos funcionales y no funcionales.}
    \label{tab:trazabilidad_requerimientos}
\end{table}

% ============================================================
\section{Conclusiones}
\label{sec:conclusiones_finales}
% ============================================================

\begin{enumerate}
    \item \textbf{La arquitectura modular permitió integrar visión y robótica
    sin acoplar la interfaz a la comunicación remota.} La separación entre el
    hilo de cámara, el worker propietario de ZeroMQ y el hilo de Tkinter evita
    que una RPC lenta congele la captura o los controles del usuario.

    \item \textbf{La selección anatómica y la asociación independiente de la
    mano resultaron esenciales para la coherencia del control.} El uso de los
    landmarks 12, 14 y 16 sobre la imagen sin espejo corrige la identificación
    del brazo derecho. La asociación de Hands por proximidad permite que la
    pérdida temporal de la mano no invalide un brazo correctamente detectado.

    \item \textbf{La validación por canales reduce la propagación de entradas
    erróneas.} Los valores no finitos, los segmentos degenerados, los comandos
    caducados y los objetivos fuera de rango se rechazan antes de actuar. Los
    buffers independientes impiden que un error del gripper contamine los
    ángulos del brazo, o viceversa.

    \item \textbf{El control del UR5/RG2 es funcional dentro del simulador y
    posee un comportamiento de fallo cerrado.} El preflight verifica la escena,
    los límites efectivos y el modo de los joints. El watchdog, el hold, la
    detección reactiva de colisiones y el ESTOP evitan continuar enviando
    objetivos cuando el estado deja de ser confiable.

    \item \textbf{Las 87 pruebas automatizadas aprobadas respaldan los
    contratos de software, no la precisión física del sistema.} Estas pruebas
    demuestran el comportamiento esperado ante entradas válidas, fallos y
    concurrencia. La integración con CoppeliaSim confirma el mapeo determinista,
    la actuación y la parada, pero no sustituye una referencia angular externa.

    \item \textbf{La simulación cuantitativa completa el análisis sin
    convertirse en evidencia experimental.} En el escenario nominal se
    obtuvieron 27.6 FPS, 63.8 ms de latencia, MAE de $3.8^\circ$ para hombro y
    $4.6^\circ$ para codo, y 94.0\% de exactitud del gripper. Los escenarios
    sintéticos predicen degradación con poca luz, movimiento rápido, oclusión y
    distancia larga; todos estos valores deben validarse con cámara.

    \item \textbf{La principal limitación proviene de la percepción monocular
    2D.} Sin profundidad métrica ni calibración, el sistema no reproduce la
    pose cartesiana completa del brazo humano. Además, el umbral del gripper en
    píxeles depende del tamaño aparente de la mano y de su distancia a la
    cámara.

    \item \textbf{El proyecto alcanza una demostración virtual reproducible,
    pero no una solución certificada para hardware real.} Las protecciones
    implementadas son adecuadas para contener fallos en CoppeliaSim. Un robot
    físico exigiría una capa de seguridad independiente, límites mecánicos,
    evaluación de riesgos y parada de emergencia por hardware.
\end{enumerate}

% ============================================================
\section{Recomendaciones}
\label{sec:recomendaciones}
% ============================================================

\begin{itemize}
    \item Ejecutar la campaña E1--E5 con la misma resolución, equipo y
    configuración, registrando cada repetición en archivos con marca temporal.
    Las tablas finales deben incluir promedio, desviación estándar, mínimo,
    máximo y número de muestras válidas, y deben reemplazar explícitamente los
    resultados sintéticos de este informe.

    \item Medir hombro y codo con un goniómetro o anotaciones calibradas que no
    dependan de MediaPipe. Solo entonces debe compararse el MAE con el objetivo
    de $5^\circ$ definido en los requerimientos.

    \item Normalizar la distancia pulgar--índice por una medida de escala de la
    mano, por ejemplo la distancia entre la muñeca y un nudillo estable. Esto
    reduciría la sensibilidad de la apertura del RG2 a la distancia de cámara.

    \item Conservar el preflight y la política de fallo cerrado al modificar la
    escena. Cualquier cambio de aliases, jerarquía, modo de joints o límites
    debe acompañarse de pruebas automatizadas y una repetición de la prueba de
    colisión externa.

    \item Fijar las versiones de Python y dependencias en un entorno
    reproducible, documentar el índice de cámara utilizado y versionar la
    escena de CoppeliaSim junto con el código que la controla.

    \item No conectar esta aplicación directamente a un UR5 físico. Antes de
    una migración deben incorporarse evaluación formal de riesgos, límites de
    velocidad y fuerza, zonas seguras, supervisión independiente y un circuito
    físico de parada de emergencia.
\end{itemize}

% ============================================================
\section{Trabajos Futuros}
\label{sec:trabajos_futuros}
% ============================================================

\begin{enumerate}
    \item \textbf{Instrumentación experimental.} Incorporar un registrador de
    FPS, latencia extremo a extremo, comandos aceptados o rechazados, eventos
    de watchdog y tiempo de respuesta del ESTOP. Los datos permitirían completar
    y sustituir las métricas simuladas sin alterar el bucle de control.

    \item \textbf{Calibración y estimación tridimensional.} Calibrar la cámara
    y evaluar landmarks 3D o una cámara de profundidad para reducir la
    dependencia del plano de imagen y estimar trayectorias cartesianas.

    \item \textbf{Mapeo cinemático completo.} Sustituir el mapeo lineal de dos
    joints por cinemática inversa con límites, control de muñeca y orientación
    del efector final. El algoritmo debe resolver la diferencia estructural
    entre el brazo humano y los seis grados de libertad del UR5.

    \item \textbf{Planificación y seguridad predictiva.} Añadir comprobación
    de autocolisiones, predicción de colisiones con el entorno y generación de
    trayectorias continuas. La actuación coordinada también podría trasladarse
    a un script de la escena para ejecutar el lote articular de forma atómica.

    \item \textbf{Adaptación automática al usuario.} Calibrar al inicio los
    rangos de movimiento y el tamaño aparente de la mano, y ajustar umbrales y
    suavizado según la distancia, la iluminación y la velocidad observadas.

    \item \textbf{Evaluación comparativa.} Comparar BlazePose con alternativas
    actuales de estimación de pose bajo el mismo protocolo, considerando
    precisión, latencia, estabilidad temporal y consumo de CPU.
\end{enumerate}

% ============================================================
\section{Cierre del Proyecto}
\label{sec:cierre_proyecto}
% ============================================================

Robot Mimic demuestra que una cámara RGB convencional y herramientas de código
abierto pueden integrarse para controlar de forma intuitiva un brazo robótico
virtual. El resultado más sólido del proyecto no es únicamente la imitación
visual, sino la incorporación de validaciones, estados y mecanismos de
contención alrededor del canal de control. La simulación cuantitativa completa
el escenario académico, pero la campaña experimental sigue siendo necesaria
para establecer el desempeño real. Con esa validación y las extensiones
propuestas, esta base puede evolucionar hacia un sistema de teleoperación más
preciso y medible sin perder la separación entre percepción, actuación y
seguridad.
